\subsection{What the controller can see}
A real junction in an Indian city has CCTV cameras, not induction loops in every lane. A camera
with a detector and a tracker (as used to build the Pune dataset) can count vehicles by class,
tell if they move, and measure how long each stopped vehicle has waited, but only within its field
of view. The agent therefore sees only the last 75\,m of each incoming lane.

\subsection{MDP formulation}
\begin{description}[leftmargin=1.6cm,style=nextline]
  \item[State] For each of the 12 incoming lanes: the PCU of halted vehicles and the PCU of all
        vehicles inside the camera view, each divided by the PCU that fits in the view; for each
        road, the longest waiting time of a vehicle in view (divided by 300\,s); a one-hot code of
        the road that has green; a flag that shows if the minimum green has passed; and the
        elapsed green time. The state has 34 values in $[0,1]$.
  \item[Action] The road that gets green next (4 actions). The agent decides every $\Delta t$
        seconds. If it keeps the same road, the green continues. If it selects another road and the
        minimum green (10\,s) has passed, the signal shows 3\,s of yellow and 2\,s of all-red and
        then the new green. A green that reaches the maximum (60\,s) passes to another road.
        In the \emph{keep/switch} variant the agent has 2 actions: keep the green, or switch to the
        road that the camera rule (below) ranks highest.
  \item[Reward] For each second $s$ of a decision step, the controller pays a cost; vehicles that
        cross the stop line earn a bonus; a change of road costs a fixed penalty:
        \begin{equation}
        r_{t+1} = -\sum_{s}\Big(w_q\,Q_s + w_b\,B_s + w_m\,M_s\Big) + w_x\,X_t - w_c\,\,\mathbf{1}[\text{switch}],
        \end{equation}
        where $Q_s$ is the PCU of stopped vehicles on the incoming lanes, $B_s$ is the PCU of
        vehicles that cannot enter the road because the queue has spilled back, $M_s$ is the
        longest accumulated waiting time of any vehicle, and $X_t$ is the PCU that crossed the stop
        line during the step. The weights are $w_q=w_b=0.01$, $w_m=0.002$, $w_x=0.2$ and
        $w_c=0.5$. The terms target queues, spill-back, starvation of a quiet road, throughput and
        the 5\,s lost at each change of road. The reward uses simulator data that a camera cannot
        give (for example $B_s$); this is allowed because the reward is used only in training,
        while the deployed controller uses only the camera state.
\end{description}

\subsection{Camera rule: a baseline and a teacher}
The camera rule is a hand-written controller that uses only camera data. It keeps the green while
vehicles are inside the last 30\,m of the green road, between the minimum and maximum green. When
the green road is clear, it serves the road with the highest score
$\sum_{\text{lanes}} \mathrm{PCU}_{\text{halted}} + 0.05\times(\text{longest wait in s})$,
so it skips empty roads. It is a baseline and also the \emph{teacher} for a warm start: before RL
training, 20{,}000 decisions of the camera rule fill the replay buffer, and the exploration rate
starts at 0.2 instead of 1.0. The agent then learns to improve on the rule.

\subsection{Strategies}
Ten training runs test one change each (Table~\ref{tab:strategies}). All use the camera state and
the combined reward unless the table says otherwise.

\begin{table}[H]
\centering\small
\begin{tabular}{lll}
\toprule
Run & Change & Question \\
\midrule
R1 & DQN, $\Delta t=5$\,s & Effect of the new reward and state alone \\
R2 & DQN, $\Delta t=2$\,s & Effect of faster decisions \\
R3--R5 & R2 + warm start from the camera rule (3 seeds) & Main candidate and its spread \\
R6 & R3 + keep/switch action & Is a 2-action design easier to learn? \\
R7 & R3 with QR-DQN & Effect of a distributional DQN \\
R8 & R3 trained for 300{,}000 steps & Effect of longer training \\
R9 & R3 with the old reward (drop in waiting time) & Contribution of the combined reward \\
R10 & R3 with a full-lane view instead of the camera & Cost of camera-only sensing \\
\bottomrule
\end{tabular}
\caption{Training runs. Each run trains for 150{,}000 steps unless noted.}
\label{tab:strategies}
\end{table}

\subsection{Training and testing}
\begin{itemize}
  \item \textbf{Training:} days 11--15 January at Jehangir Chowk. Each episode is one random hour
        of the morning peak with the demand scaled by a random factor in $[0.7,1.4]$.
        DQN from Stable-Baselines3, MLP with two hidden layers of 256 units,
        learning rate $5\times10^{-4}$, replay buffer $2\times10^5$, batch 128, $\gamma=0.99$,
        target update every 2{,}000 steps. Each run uses 4 parallel SUMO instances on a Kaggle
        notebook.
  \item \textbf{Test scenarios:} days 16--18 January, full 3-hour peak, three route seeds per day:
        (1) unseen days at Jehangir Chowk; (2) the same days at 1.3 times the real demand;
        (3) an unseen junction, RTO Chowk; (4--5) unseen days at 1.0 and 1.3 times demand with
        free left turns. Each route seed fixes the vehicle departures, so all controllers see the
        same vehicles.
  \item \textbf{Baselines:} a fixed-time plan (30\,s green per road, 140\,s cycle); the SUMO
        actuated controller on the same four phases, which uses an induction loop on every lane
        and extends a green between 10 and 60\,s while vehicles arrive; and the camera rule.
  \item \textbf{Metrics:} from SUMO trip information. Waiting time and delay include the time a
        vehicle waits to enter the network. Throughput is the number of vehicles that complete
        their trip per hour. Fuel and CO$_2$ are per vehicle.
\end{itemize}

